% Replacement for Sections 5.2 and 5.3 (main.tex l.732--901).
% Draft 2026-11-11, written against the measured autoregressive ladder.
%
% The section this replaces argued that the causal contribution falls as the
% evaluator saturates.  Measured across \TrendNEval\ autoregressive
% evaluators spanning \TrendCapSpan$\times$ in parameters, it does not.  What
% rises is the fraction of texts already at the ceiling.  The two are easy to
% confuse and the earlier draft confused them; the rank correlations below
% are reported precisely so the text cannot.

\subsection{What the measured distance is made of}
\label{sec:contrast}

Section~\ref{sec:blkord} supplies the origin that Eq.~\eqref{eq:delta}
lacks: the distribution of $\Delta$ over the $n!$ arrangements of a text's
own sentences, against which its authored arrangement has a position.
Section~\ref{sec:decomp} splits the median position's distance from the
$0.50$ null into the part a non-causal --- but still authored --- order
earns, and the part causal order earns beyond it.

Table~\ref{tab:decomposition}, introduced with the rule in Section~\ref{sec:decomp}, gives both for every evaluator.  The authored
component is $\AuthoredGPTtwo$ at $\ParamsGPTtwo$ and $\AuthoredQwen$ at
$\ParamsQwen$; the causal component is $\CausalGPTtwo$ and $\CausalQwen$.
Between a tenth and a fifth of the measured distance is attributable to the
labelled contrast, and the rest to the arrangement simply being the one a
writer chose.

\subsection{The causal component is one step of the statistic}
\label{sec:resolution}

Before reading those numbers as magnitudes, note what the statistic can
express.  The midrank over $n!$ arrangements lives on a grid of width
$1/n! = \GridStep$.  A causal component of $\CausalGPTtwo$ is one grid step.
Across the \CausalCompNEval\ autoregressive evaluators the component takes
\CausalCompNDistinct\ distinct values, between \CausalCompStepsMin\ and
\CausalCompStepsMax\ steps, while the authored component spans several.

The causal contribution is therefore at the instrument's resolution.  Its
value is not a measured size but the smallest non-zero reading available,
and a difference between two evaluators of one step is a difference between
adjacent representable numbers.  This is why Section~\ref{sec:decomp}
reports the same split on
$z = \Delta(\mathrm{id})/\mathrm{sd}\{\Delta(\pi)\}$, which has neither a
ceiling nor a grid: there the share runs $\ShareZMin\%$ to $\ShareZMax\%$.

Nothing here says the causal contribution is absent.  The rank test of
Section~\ref{sec:blkord} separates the two label groups at
$p < \BlkPMaxCausal$ in every corpus and evaluator.  What is at resolution
is the \emph{size} of that separation once it is expressed as a share of the
authored distance, and only the $z$ form can carry it.

\subsection{The room runs out; the contribution does not fall}
\label{sec:window}

A strong evaluator pushes an arbitrary authored order towards the top of its
own permutation range, and once the authored order \emph{is} the extreme
arrangement no normalisation can express a further increase.  That ceiling
is real and it is approaching.  Over the \TrendCapSpan$\times$ span of
parameters covered here, the fraction of causally labelled texts whose
authored order is already the extreme arrangement rises with capacity at
$\rho = \TrendCeilRho$ ($p = \TrendCeilP$, Spearman, $n = \TrendNEval$) ---
the only component with a significant trend.

The causal component itself does not fall.  Against capacity it gives
$\rho = \TrendCausalRho$ ($p = \TrendCausalP$), and the ceiling-free $z$
share gives $\rho = \TrendShareZRho$ ($p = \TrendShareZP$).  Both
correlations are taken after rounding the components to $10^{-\TrendQuantum}$,
because at this resolution they are differences of medians on the same grid:
the causal component takes \TrendCausalNDistinct\ distinct values across the
\TrendNEval\ evaluators, and without rounding, differences of order
$10^{-16}$ between values that are mathematically equal would be ranked as
an ordering.  The authored
component rises, $\rho = \TrendAuthoredRho$ ($p = \TrendAuthoredP$), which
is the saturation showing up directly.

The distinction matters because the two readings make different predictions.
``The causal contribution shrinks with capacity'' would say the quantity
responds less and less to causal order, and extrapolates to a measurement
that eventually sees nothing.  What is measured is the opposite in kind: the
causal contribution is constant to within one grid step across
\TrendCapSpan$\times$ in capacity, while the space above an authored order
closes.  The instrument loses headroom, not signal --- and the signal was at
the resolution floor from the smallest evaluator onwards.

We therefore do not extrapolate past $\ParamsQwen$, and we claim only what
the ladder supports: a tenth to a fifth of what $\Delta$ measures is
attributable to causal order, that share does not vary systematically with
capacity, and the margin within which it could be measured at all is
narrowing.
